\section{Lattice VOA, canonical involution and order-two Moonshine branch}
\label{sec:voa-moonshine-orden-dos-rev11}
\index{VOA!Leech lattice}
\index{Moonshine!class 2A}

The incidence branch first constructs the Leech lattice
\(\Lambda_{24}\). From that datum, the classical theory of vertex operator
algebras provides a canonical prolongation. If
\(\mathfrak h=\Lambda_{24}\otimes_{\mathbb Z}\mathbb C\) and \(M(1)\) is
the Fock module of its Heisenberg algebra, the lattice VOA is
\begin{equation}
 V_{\Lambda}=M(1)\otimes\mathbb C[\Lambda_{24}],
 \label{eq:voa-reticular-leech-rev11}
\end{equation}
with central charge \(24\). For \(\alpha\in\Lambda_{24}\), its vertex
operators are written
\[
 Y(e^\alpha,z)=E^-(-\alpha,z)E^+(-\alpha,z)e^\alpha z^{\alpha_0}.
\]
The absence of roots in \(\Lambda_{24}\) implies the vanishing of the
weight-one sector in the Moonshine orbifold.

The involution
\[
 \theta:\Lambda_{24}\longrightarrow\Lambda_{24},
 \qquad \theta(\lambda)=-\lambda,
\]
is intrinsic and prolongs to \(V_{\Lambda}\). Its graded trace is obtained
by direct counting of the oscillator modes.

\begin{theorem}[Trace of the negation involution]
\label{thm:traza-theta-leech-rev11}
\begin{equation}
 t_2(\tau)
 :=\operatorname{Tr}_{V_\Lambda}
 \bigl(\theta q^{L_0-1}\bigr)
 =q^{-1}\prod_{n\geq1}(1+q^n)^{-24}.
 \label{eq:t2-producto-leech-rev11}
\end{equation}
\end{theorem}
\begin{proof}
The involution pairs \(e^\lambda\) with \(e^{-\lambda}\); since the lattice
is torsion-free, only \(\lambda=0\) contributes to the trace of the
lattice sector. Each one of the twenty-four energy oscillators \(n\) contributes
\(\sum_{k\geq0}(-1)^kq^{nk}=(1+q^n)^{-1}\). The vacuum displacement
produces the factor \(q^{-1}\).
\end{proof}

The Fricke completion associated with this trace is
\begin{equation}
 T_{2A}(\tau)=t_2(\tau)+24+\frac{2^{12}}{t_2(\tau)}
 =q^{-1}+4372q+96256q^2+1240002q^3+\cdots.
 \label{eq:T2A-hmt-rev11}
\end{equation}
The term \(24\) cancels the constant term of \(t_2\); the factor
\(2^{12}\) is the degeneracy of the twisted sector of twenty-four bosons.
The same boundary series determines the modular invariant through
\begin{equation}
 \boxed{j(\tau)=\frac{(t_2(\tau)+256)^3}{t_2(\tau)^2}}.
 \label{eq:j-desde-t2-rev11}
\end{equation}
Equations \eqref{eq:t2-producto-leech-rev11}--
\eqref{eq:j-desde-t2-rev11} construct the canonical branch of order two without
continuous parameters.

The Frenkel--Lepowsky--Meurman orbifold is defined by
\[
 V^\natural=(V_\Lambda)^+\oplus(V_\Lambda^T)^+.
\]
It is a holomorphic VOA of central charge \(24\), satisfies
\((V^\natural)_1=0\), and has character
\[
 \operatorname{ch}_{V^\natural}(\tau)=j(\tau)-744
\]
and its automorphism group is the Monster group
\cite{FLM1988,ConwayNorton1979,Borcherds1992}. In this chain, HMT fixes the
Leech lattice and the involution \(\theta\); the classical VOA and
FLM theorems identify the resulting prolongation with \(V^\natural\) and its
automorphism group. The assertion does not require an action of the Monster on the
digits: the Archimedean evaluation and the exceptional branch descend from a
common HMT preimage and preserve different invariants.

The neighboring function
\(T_{3A}=t_3+12+729/t_3\) belongs to the classical theory of Hauptmoduln of
order three. The active HMT realization in the tritic corridor is the class
\(3B\); an identification with \(3A\) requires declaration of the specific lift
that changes the class.
